\section{Long Autonomous Rollouts}
\label{app:long-horizon}

\paragraph{Test case and scope.}
We extend the selected seed-zero maps to 4,096 successive forecast intervals from one held-out initial history: trajectory 180, frames 0--4. Each model carries eight independent generated histories after initialization. The It\^o map uses eight generative steps per frame and the flow map one; matched 64-step stochastic-interpolant and flow-matching references use the same initial history and ensemble size. This extension is separate from the all-20-trajectory, three-student confirmation in the main text.

\paragraph{Time and reference data.}
The recovered companion archive supplies a mean saved-time increment of 0.49989997016; thus this horizon spans approximately 2,047.59 nominal simulation-time units. The primary coordinate remains stored snapshot intervals. Exact matching of all coarse fields to the companion's stride-two downsampling establishes the timestamp association, but does not identify the original viscosity or forcing parameters. Only 95 future truth fields are available, corresponding to archive frames 5--99. RMSE and spread/skill stop at lead 95; later diagnostics do not measure accuracy against a reference continuation.

\paragraph{Diagnostics.}
We inverse-transform each member into the archive's vorticity units, without changing the internal model trajectory. For $P$ spatial cells, member $e$ has area-mean enstrophy $Z_e=(2P)^{-1}\sum_p\omega_{ep}^2$. We report its ensemble mean, rather than the generally smaller enstrophy of the ensemble-mean field. In fact,
\begin{equation}
\frac1E\sum_e Z_e
= Z(\bar\omega)+\frac{E-1}{2E}S_{\mathrm{RMS}}^2,
\end{equation}
where $S_{\mathrm{RMS}}$ uses unbiased sample variance as in Appendix~\ref{app:metrics}. Corrected spread is $\sqrt{(E+1)/E}\,S_{\mathrm{RMS}}$; the ratio divides it by ensemble-mean RMSE. No amplitude clipping, rescaling, truth reset or replacement of a failed member is allowed.

\begin{table*}[t]
\caption{Long autonomous rollouts from one history with eight members. Late statistics average the final 512 leads; $Q_{1:95}$ is the ratio of pooled corrected RMS spread to pooled RMSE over the available-reference interval.}
\label{tab:long-maps}
\centering
\pendingbox{rows: It\^o map (8 steps), flow map (1 step), SI64, FM64, and the stored truth over leads 1--95; columns: completed leads, finiteness, mean member enstrophy over leads 1--95 and over the final 512 leads, pooled RMSE and $Q_{1:95}$ over leads 1--95.}
\end{table*}

\paragraph{Observed behavior.}
\pending{long-horizon behavior: boundedness of all four trajectories, retained member enstrophy and spread relative to the truth over leads 1--95 and to the contextual archive mean, and late-lead amplitude of each model (Table~\ref{tab:long-maps}, Figure~\ref{fig:long-maps-skill}).} These correlated leads from one initial condition do not provide an independent-case calibration estimate, and boundedness of these runs does not establish accurate invariant statistics or stability for other initial conditions.

\paragraph{Spatial enstrophy spectra.}
To resolve the scales of the first and last generated fields, let
$\widehat\omega_e(q)=P^{-1}\sum_p\omega_{ep}\exp(-2\pi i\,q\cdot p/n)$
for the $n\times n$ grid, and group integer wavevectors by nearest radial shell. The shell-summed enstrophy is
\begin{equation}
\mathcal Z_e(k)=\frac12\sum_{q:\,\lfloor\lVert q\rVert+1/2\rfloor=k}
|\widehat\omega_e(q)|^2.
\end{equation}
Parseval's identity gives $\sum_{k\ge0}\mathcal Z_e(k)=Z_e$.
This is a sum over modes in each shell, not a radially averaged power spectrum; it does not redefine the earlier spectral-error metric. The spectra describe the observed $128^2$ grid without additional filtering. Frequencies are cycles per grid domain; no physical domain length is assumed. Figures~\ref{fig:long-spectrum-first} and~\ref{fig:long-spectrum-last} show separate first- and last-step spectra on log--log axes. The zero-frequency component is omitted from these logarithmic axes. All square-grid modes are included; shells above the axis Nyquist frequency $k=64$ contain only corner modes. No spectral values are floored. Reference truth is available for the first forecast step only. \pending{shell-by-shell comparison of the models at lead 4,096.}

\begin{figure*}[t]
\centering
\pendingbox{first-step (lead 1) shell-summed enstrophy spectra: member-mean curves with 5--95\% member bands for each model, and the matching observed field.}
\caption{Spatial enstrophy spectrum at the first forecast step (lead 1). Curves are means of eight member shell-summed spectra, with member 5--95\% bands, in archive vorticity units; the observed field is shown for reference. Both axes are logarithmic; $k$ denotes radial spatial frequency in cycles per grid domain, and higher shells than the axis Nyquist frequency 64 contain square-grid corner modes.}
\label{fig:long-spectrum-first}
\end{figure*}

\begin{figure*}[t]
\centering
\pendingbox{last-step (lead 4,096) shell-summed enstrophy spectra for each model, same normalization and axes as Figure~\ref{fig:long-spectrum-first}.}
\caption{Spatial enstrophy spectrum at the last forecast step (lead 4,096), using the same normalization and log--log spatial-frequency axes as Figure~\ref{fig:long-spectrum-first}. Member-mean spectra and 5--95\% member bands describe generated fields only: no reference continuation is available at this lead. Differences between models therefore show different retained spatial variability, not verified long-time spectral accuracy.}
\label{fig:long-spectrum-last}
\end{figure*}

\begin{figure*}[t]
\centering
\pendingbox{available-reference diagnostics over leads 1--95: corrected spread and ensemble-mean RMSE per lead (left) and their ratio (right), for each model.}
\caption{Available-reference diagnostics for the long rollouts, restricted to leads 1--95. The correction is $\sqrt{9/8}$ applied to unbiased RMS spread for eight members. Similar RMSE does not compensate for smaller spread. No skill values are extrapolated beyond the stored reference.}
\label{fig:long-maps-skill}
\end{figure*}
